\chapter{What restricted access cannot upper-bound}
\label{ch:dissipation-limits}

Sharper lower certificates do not automatically make total dissipation
identifiable. This chapter separates three obstructions: an exactly hidden
current, an identifiable parameter with a singular thermodynamic dependence,
and incomplete physical channel resolution. Its main construction uses one
connected four-state cycle. Every microscopic jump changes the detector output.
There is no independent hidden subsystem, no parallel microscopic edge, and no
unobserved transition time.

Strong lumpability and loss of entropy production under coarse-graining are
established results~\cite{ac_buchholz,ciBusiello2019}. Optimizing over compatible
hidden mechanisms is also an established approach to sharp lower
bounds~\cite{ciEhrich2021,ciNitzan2023}. The contribution developed here is the
explicit obstruction-and-repair chain: complete observed path equivalence at
fixed dimension and activity, a quantitative finite-exposure impossibility that
survives full microscopic observation, and calibrated routes to usable upper
intervals. These proofs do not establish historical priority or a new physical
law.

\section{Four states, visible jumps, unbounded hidden dissipation}
Number the states $0,1,2,3$ around a simple cycle. Let $r_i>0$ be known exit
rates and $0<q\le1/2$ an unknown reverse probability. Set
\begin{equation}
 Q_{i,i+1}=r_i(1-q),\qquad Q_{i,i-1}=r_iq,\qquad Q_{ii}=-r_i,
                                                        \label{eq:dl-cycle}
\end{equation}
with indices modulo four. The stationary law, mean jump activity, and total
stationary entropy-production rate are
\begin{align}
 \pi_i&=\frac{r_i^{-1}}{\sum_j r_j^{-1}},\qquad
 K=\frac4{\sum_j r_j^{-1}},\\
 \sigma(q)&=K(1-2q)\log\frac{1-q}{q}.                    \label{eq:dl-entropy}
\end{align}
Indeed the stationary clockwise and counterclockwise fluxes across each
undirected edge are respectively $(1-q)/\sum_jr_j^{-1}$ and
$q/\sum_jr_j^{-1}$. Summing their current times flux affinity gives
\eqref{eq:dl-entropy}. The chain is irreducible for every allowed $q$, and all
four states participate in the observed dynamics.

Let the detector report only the parity $Z=i\bmod2$. Every microscopic jump
flips $Z$. If all $r_i=s$, the full observed process is exactly the symmetric
two-state chain with jump rate $s$, independently of $q$. Its waiting times are
independent exponentials of rate $s$, including the entire timed jump record.
At stationarity its binary autocorrelation is $e^{-2st}$ and its zero-lag
amplitude is exactly one. Thus even perfect calibration of that amplitude,
activity, state count, and graph does not distinguish the models.

Meanwhile $\sigma(1/2)=0$, and $\sigma(q)\to\infty$ as $q\downarrow0$.
Every intermediate nonnegative entropy rate occurs. The obstruction persists
under arbitrary feedback that multiplies both outgoing rates by a bounded
factor chosen from the observed parity history: the observed jump intensity
remains $s$ times that factor and still contains no $q$.

This is an ordinary lumping obstruction, made explicit in a case where all
jumps are visible. It does not say that the microscopic rates are all known:
the exit rates are known, while their division between clockwise and
counterclockwise motion is not. A known exit-rate floor is not a known
reverse-transition floor.

\section{Breaking exact blindness by kinetic tags}
Distinct exit rates can encode microscopic state information in the waiting
times without changing the binary detector. Write
$f_i(t)=r_i e^{-r_i t}$. At a typical stationary jump into even parity, the
microscopic entry state is uniform on $\{0,2\}$ because stationary jump fluxes
are uniform. The next two holding times have joint density
\begin{equation}
 g_q(u,v)=\frac{(f_0+f_2)(u)(f_1+f_3)(v)}4
       +\frac{1-2q}{4}(f_0-f_2)(u)(f_1-f_3)(v).
                                                        \label{eq:dl-pair}
\end{equation}
Consequently
\begin{equation}
 \operatorname{TV}(g_q,g_{q'})
 =|q-q'|\operatorname{TV}(\operatorname{Exp}r_0,\operatorname{Exp}r_2)
          \operatorname{TV}(\operatorname{Exp}r_1,\operatorname{Exp}r_3).
                                                        \label{eq:dl-pair-tv}
\end{equation}
This follows directly by subtracting the densities and factoring the absolute
integral. The bias is therefore identifiable within this calibrated four-state
family whenever both parity classes have distinct exit rates.

The distribution in \eqref{eq:dl-pair} is a Palm distribution: a typical event
sampled from the stationary jump process. It need not be the distribution of
the first even entry after an arbitrary clock-time origin. Successive pairs
in one trajectory are dependent. Neither event selection nor independence is
silently granted by this identity.

The repair has a quantitative weakness. For tags
\[
 (r_0,r_1,r_2,r_3)=(1+\eta,1+2\eta,1-\eta,1-2\eta),
 \qquad0<\eta<1/2,
\]
both exponential-law distances are of order $\eta$, so pair-law separation is
of order $\eta^2$. For the event $u\le1,v\le1$, the absolute derivative of its
probability with respect to $q$ is exactly
\begin{equation}
 2e^{-2}\sinh\eta\sinh(2\eta)
   =4e^{-2}\eta^2+O(\eta^4).                            \label{eq:dl-tag-signal}
\end{equation}
An independently sampled binary-window experiment based on that difference
would have a conservative Hoeffding record requirement of order $\eta^{-4}$
at fixed parameter separation. Obtaining independent Palm pairs is an additional
experimental requirement. No free reset or independent sampling assumption is
included in this scaling statement.

More generally, the conditional parity jump intensity always lies between the
two exit rates of the current parity class. For any two biases with the same
$r_i$, the entire coarse record over total exposure at most $L$ obeys
\begin{equation}
 \operatorname{TV}(P_q^Z,P_{q'}^Z)
 \le\min\{1,L\max(|r_0-r_2|,|r_1-r_3|)\}.              \label{eq:dl-coarse-coupling}
\end{equation}
Couple the two predictable jump intensities while their observed histories
agree. Their difference is bounded by the displayed rate spread; the probability
of any unmatched jump is at most its integrated bound. Initial parity laws
agree because $\pi$ does not depend on $q$. This shows explicitly how nearly
lumpable connected models can remain close over a fixed experimental budget.

\section{Full microscopic access does not remove a finite-budget singularity}
Even after kinetic tags make the parameter identifiable, the entropy function
is singular near a vanishing reverse rate. The following result concerns the
\emph{complete microscopic record}, so every coarser observation inherits it.

\begin{theorem}[Finite-exposure dissipation ambiguity]
\label{thm:dl-budget}
Fix any four positive exit rates $r_i$, a total pathwise exposure budget $L>0$,
a tolerance $0<\varepsilon<1$, and any desired entropy gap $M>0$. There are two
strictly positive reverse probabilities $q_0,q_1$ in \eqref{eq:dl-cycle}, with
the same state count, graph, stationary law, exit rates, and activity, such that
\[
 |\sigma(q_1)-\sigma(q_0)|>M,
 \qquad \operatorname{TV}(P_{q_0},P_{q_1})\le\varepsilon.
\]
The record may contain all microscopic states, directions and jump times,
including multiple independent resets with the common initial law and total
exposure at most $L$.
\end{theorem}
\begin{proof}
Use one common uniformization clock of rate $R=\max_i r_i$. While the two
microscopic paths agree, their holding and acceptance decisions can be identical.
An accepted jump has different direction with probability $|q_0-q_1|$ under
a common uniform direction draw. Hence
\begin{equation}
 \operatorname{TV}(P_q,P_{q'})\le\min\{1,RL|q-q'|\}.
                                                        \label{eq:dl-full-coupling}
\end{equation}
Take $q_0=\min\{1/4,\varepsilon/(2RL)\}$ and then let $q_1>0$ approach zero.
The record bound remains at most $\varepsilon/2$, whereas
$\sigma(q_1)\to\infty$. All chosen finite models retain strictly positive
reverse rates. Bounded controls that preserve this common direction parameter
and a common total exit-rate cap admit the same coupling argument.
\end{proof}

A standard two-point testing argument now gives a minimax consequence. For any
estimator and either member of a pair separated by $M$, at least one error
probability for absolute error at least $M/2$ is at least
$(1-\varepsilon)/2$. Finite dimension, a known graph, bounded activity and a
positive exit-rate floor do not provide a uniform additive-error estimate of
total dissipation.

The exact executable example fixes
\[
 (r_0,r_1,r_2,r_3)=(1.02,1.04,0.98,0.96),\qquad L=10.
\]
Two finite models have a certified entropy-rate gap exceeding $209.8058$
while their full-record total variation is at most $0.005$. The units are
natural entropy per declared model time. Large entropy differences arise from
an uncalibrated large affinity, not from a larger state space, faster activity,
or an appended hidden subsystem.

\section{Why a valid upper confidence bound must sometimes be infinite}
\begin{theorem}[Unavoidable unbounded upper intervals]
Suppose $U$ is an extended-real upper confidence bound computed from an allowed
record with exposure at most $L$, and
$\Pr_q[U\ge\sigma(q)]\ge1-\alpha$ for every $0<q\le1/2$.
Then, for each such finite model,
\begin{equation}
 \Pr_q[U=\infty]\ge\max\{0,1-\alpha-RLq\}.            \label{eq:dl-infinity}
\end{equation}
\end{theorem}
\begin{proof}
Introduce only as a mathematical limit the clockwise process at $q=0$.
Equation~\eqref{eq:dl-full-coupling} gives
$\operatorname{TV}(P_q,P_0)\le RLq$. Along $q_j\downarrow0$, coverage implies
$\Pr_0[U\ge\sigma(q_j)]\ge1-\alpha-RLq_j$. Since
$\sigma(q_j)\uparrow\infty$, continuity from above gives
$\Pr_0[U=\infty]\ge1-\alpha$. Transfer this event back to a finite $q$ with
the same total-variation bound.
\end{proof}

For the finite reference model in the example, every uniformly valid $95\%$
upper-bound procedure must return infinity with probability at least $94.5\%$.
This does \emph{not} mean that every data-dependent upper bound must always be
infinite. Enough resolved reverse events can exclude very small $q$ and produce
a finite upper interval. The obstruction is to uniformly finite guarantees
without further restrictions on the admissible reverse rates.

\section{Small access repairs and calibrated physical caps}
One additional binary tag identifying a single physical edge already repairs
structural identifiability within this ring. The existing parity before and
after that tagged edge determines its direction. Its stationary fluxes
$F_+,F_-$ give
\begin{equation}
 \sigma=4(F_+-F_-)\log(F_+/F_-).                        \label{eq:dl-one-edge}
\end{equation}
Every edge has the same stationary cycle current and flux ratio. This is an
instance of established unicyclic inference from a resolved edge, closely related
to visible-transition methods~\cite{ciHarunari2022}. Finite-time edge events
remain dependent, and \eqref{eq:dl-one-edge} does not itself supply an independent
binomial sampling law.

A calibrated affinity cap provides a different repair. If every stationary
undirected edge flux pair $F_e,B_e$ obeys
$|\log(F_e/B_e)|\le A$, then
\begin{equation}
 \sigma\le K A\tanh(A/2).                              \label{eq:dl-affinity-cap}
\end{equation}
Indeed its contribution is
$(F_e+B_e)a_e\tanh(a_e/2)$, where $a_e=\log(F_e/B_e)$, and this even function
increases with $|a_e|$. Sum over edges. The four-cycle attains equality when
all edge flux affinities have magnitude $A$, so this bound is sharp under the
stated information.

A stationary \emph{flux} affinity includes the stationary population ratio;
it is not automatically a calibrated bound on $\log(Q_{ij}/Q_{ji})$.
If instead the latter bare rate ratios have a known bound $B$, stationarity
cancels the population-potential term and gives the safe, generally looser
bound $\sigma\le KB$. In particular, known bounds on \emph{every individual
forward and reverse rate}, $m\le Q_{ij}\le M$, imply
$\sigma\le K\log(M/m)$. Merely knowing the exit rates $r_i$ does not supply
such a reverse-rate floor.

\section{A finite confidence interval with declared observation time}
For an executable finite-count repair, suppose the direction of every jump is
resolved. Collect a predetermined quota of $n$ jumps. Their reverse indicators
are independent Bernoulli($q$) draws, even with unequal exit rates. This statement
must not be transferred to the number of directions observed by a fixed time:
conditioning on that random jump count can bias directions because rates depend
on the visited states.

Exact binomial inversion produces simultaneous bounds $q\in[q_-,q_+]$.
Intersect these with any independently calibrated floor $q\ge\delta$ and with
$q\le1/2$, giving $q_L=\max\{q_-,\delta\}$ and
$q_U=\min\{q_+,1/2\}$. An empty intersection flags incompatibility of the
record and the declared parameter class. Otherwise, since
\eqref{eq:dl-entropy} decreases on $(0,1/2]$, the entropy interval is
\[
 [\sigma(q_U),\sigma(q_L)].
\]
If the lower endpoint for $q$ is zero, the upper entropy endpoint is infinity.
With a positive floor it is finite for every dataset compatible with the model;
without a floor, positive reverse evidence can still make it finite.

For the tagged rates above, $n=1000$, ten reverse observations, and a calibrated
floor $\delta=0.001$, exact rational binomial and logarithm bounds give a $95\%$
interval contained in
\[
 [3.83196,\;5.27668]
\]
natural entropy units per model time. The prior floor alone only supplies the
upper bound $6.88605$. These counts are an explicitly synthetic worked input.
With the same positive reverse evidence the interval is also finite without a
prior floor; with zero reverse counts and no floor it is unbounded above.

Observation duration is a separate resource. If $r_{\min}=\min_i r_i$, the
$n$ holding times are dominated by independent exponentials of rate
$r_{\min}$, so $\mathbb E T_n\le n/r_{\min}$. A timeout at
$2n/r_{\min}$ has failure probability at most
\[
 (2/e)^n\le(3/4)^n.
\]
For this example the expected-time bound is $1041.67$, the timeout is
$2083.33$, and the timeout bound is below $1.16\times10^{-125}$. A timeout
returns the full prior interval. Selecting only completed runs and calling the
result an unbiased conditional binomial interval is not justified. Preparation,
reset, detector and control overheads remain additional unmeasured costs.

\section{Scope and reproducibility}
The exact generator identities, entropy intervals, finite-budget pairs,
count inversion, and timing bounds are reproduced by
\texttt{dissipation\_limits.py}. Frozen-seed uniformization simulations check
complete and parity record couplings. Their mismatch frequencies are properties
of the chosen coupling, not estimates of total variation. Thirteen producer
tests check exact stationarity, full parity-path likelihoods, independent
waiting-density integrals, weak-tag scaling, finite-budget ambiguity, and
resolved-count inference.

The thermodynamic interpretation concerns a specified finite Markov jump
network with physically resolved reverse transitions and the appropriate local
detailed-balance accounting. A compatible embedding takes state-energy
differences proportional to $\log r_i$ and a uniform clockwise driving work
proportional to $\log[(1-q)/q]$. If additional channels or unresolved physical
states exist, they introduce further ambiguities already discussed elsewhere
in the reference. The present obstruction occurs before those complications
are added.
